\subsection{Throughput guards, and the result the gridlock was hiding}\label{sec:res-guards}
\label{sec:res-calib}
Mean delay can be improved by stranding vehicles, so \S\ref{sec:met-metrics}
commits to guard metrics alongside it. Table~\ref{tab:guards} changes how
two of the three topologies read.

\begin{table*}[t]
\centering
\small
\begin{tabular}{llrrrr}
\toprule
Topology & Arm & Loaded & Compl. & Teleports & Undep. \\
\midrule
Euston peak-hour & fixed-time & 780 & 51.1\% & 22.2 & 184 \\
 & MaxPressure & 780 & 45.3\% & 31.1 & 218 \\
 & qwen ft & 780 & 48.7\% & 28.1 & 200 \\
 & qwen stock & 780 & 50.3\% & 23.3 & 199 \\
Bloomsbury grid & fixed-time & 1715 & 15.1\% & 115.0 & 328 \\
 & MaxPressure & 1715 & 14.6\% & 119.3 & 336 \\
 & qwen ft & 1715 & 14.2\% & 120.0 & 341 \\
 & qwen stock & 1715 & 13.5\% & 135.6 & 348 \\
Old Street & fixed-time & 480 & 83.3\% & 0.7 & 1 \\
 & MaxPressure & 480 & 85.8\% & 1.3 & 0 \\
 & qwen ft & 480 & 86.7\% & 0.5 & 0 \\
 & qwen stock & 480 & 84.4\% & 1.0 & 1 \\
\bottomrule
\end{tabular}
\caption{Guard metrics over the same $n{=}30$ training-disjoint seeds as
every delay comparison. Compl.\ is completed vehicles as a share of those
loaded within the 1{,}200-step horizon; teleports are SUMO's
gridlock-escape removals.}
\label{tab:guards}
\end{table*}

\finding{Of the three original cells, only Old Street is measured in a
free-flowing regime.} It clears, teleports are effectively absent, mean
delay is well-posed, and every arm significantly beats the fixed-time floor
there (Table~\ref{tab:robust}). Euston sits between, at 45--51\%
completion. There the SLM arms \emph{complete more vehicles than
MaxPressure} (380 and 392 against 353) without a significant delay
advantage, though not more than the fixed-time floor's 399; no
significance test was run on completion. \risk{Uncalibrated Bloomsbury is
a gridlock-regime measurement.} Under 15\% complete and over a hundred
teleport out of deadlock, so its mean delay compares two congested-failure
modes. The stock arm completes fewer vehicles and teleports more than
MaxPressure, so the stock harm survives the guard test in direction, but it
is regime-scoped. Completion also reproduces \S\ref{sec:res-maxpressure}'s
ordering without using delay, the tie appearing as MaxPressure marginally
worse (14.6\% against 15.1\%). So \finding{the sign flip is not an
artefact of a fragile mean}, though completion is as consistent with
saturation or link storage as with geometry.

Bloomsbury carried the study's only significant closed-loop result against
MaxPressure, so I re-examined it. Tripling the horizon to 3{,}600 steps
raised completion only to 35\% while multiplying teleports roughly
sevenfold (121 to 834): the network was over capacity, not
under-simulated. Subsampled demand located the boundary sharply, clearing
at 40\% (87.6\% completion, zero teleports) and collapsing again at 55\%
(57.5\%), so I added a demand-calibrated variant at 40\% (686 vehicles)
and re-ran all three arms over the identical thirty seeds, retaining the
oversaturated cells. Whether this is a principled correction or a search
for a more favourable test bed is argued in
Appendix~\ref{app:calib-audit}: the guard metrics predate the decision to
look, the 40\% target matched Old Street's completion rate rather than any
effect size, and the rule was applied uniformly to Euston peak-hour, which
also fails the guard and is where SLM control looks \emph{worst}
(\S\ref{sec:res-factorial}), so calibrating it risked the conclusion rather
than protecting it.

\finding{Calibrating Euston does not rescue the corridor, and the
distillation effect does not survive there at full size.} Completion rises
from 45--51\% to 68--74\%, but in the MaxPressure arm 7.7 vehicles per run
still teleport and 56.9 never depart, whereas on the calibrated grid both
classical comparators record 0.0 teleports and 0.0 undeparted and only the
stock arm differs (1.70 teleports, 5.7 undeparted). On the partly cleared
corridor the fine-tuned arm is still $+8.5\%$ against the fixed-time floor
($[-2.4, +19.5]$, $p{=}.12$) and the stock arm \risk{$+34.8\%$}
($[+5.3, +64.4]$, $p{=}.023$). The mean spares the fine-tuned arm only
because the tail is enormous: it beats the floor on six of thirty seeds,
its median seed is $+11.7\%$, and the rank test rejects ($p{=}.016$).
Neither arm differs significantly from MaxPressure ($-0.3\%$ at $p{=}.93$,
$+21.9\%$ at $p{=}.062$). The head-to-head favours fine-tuning in direction
and on the typical seed (median $-7.4\%$, 23 of 30 seeds, Wilcoxon
$p{=}.045$), but its mean is $-3.7\%$ and not significant
($[-19.1, +11.7]$, $p{=}.63$), across per-seed outcomes from $-71.6\%$ to
$+104.5\%$.

Two conclusions follow, the first against me: \finding{saturation explains
the Bloomsbury result; on Euston it is reduced, not tested}. The same
correction yields on the corridor a fifth of the grid's magnitude and no
significant mean, though in absolute delay ($-27.4$\,s against
$-35.9$\,s) the two are comparable. I cannot attribute the surviving
deficit to topology, since nothing here varies geometry while holding
demand and storage fixed. Topology and residual saturation both stay live,
and the grid effect did not replicate on the second network I calibrated.
The other conclusion concerns method. The threshold was applied where it
could damage the conclusion, and did: the rule, take the largest demand
fraction that still clears, produced an unfavourable result, reported as
such.

\begin{table*}[t]
\centering
\small
\begin{tabular}{llrrrr}
\toprule
Regime & Arm & $\Delta$ vs fixed-time & $\Delta$ vs MaxPressure & Compl. & Auth. \\
\midrule
Oversaturated & qwen ft & $+0.8$ $[-1.0, +2.6]$ & $+0.7$ $[-0.7, +2.1]$ & 14.2\% & 0.70 \\
(14\% clears) & qwen stock & $\mathbf{+2.8}$ $[+0.9, +4.6]$ & $\mathbf{+2.6}$ $[+1.1, +4.1]$ & 13.5\% & 0.62 \\
\midrule
Calibrated & qwen ft & $\mathbf{-8.8}$ $[-9.4, -8.2]$ & $+0.00$ $[-0.56, +0.56]$ & 87.4\% & 0.76 \\
(87\% clears) & phi ft & $\mathbf{-8.8}$ $[-9.3, -8.3]$ & $+0.03$ $[-0.43, +0.48]$ & 87.5\% & 0.77 \\
 & qwen stock & $\mathbf{+11.9}$ $[+10.0, +13.8]$ & $\mathbf{+22.7}$ $[+20.6, +24.7]$ & 76.4\% & 0.59 \\
\bottomrule
\end{tabular}
\caption{The same controllers under two demand regimes on the same
Bloomsbury network ($n{=}30$ paired seeds each); positive $\Delta$ is worse
than the fixed-time floor.}
\label{tab:calib}
\end{table*}

\begin{figure*}[t]
\centering
\begin{tikzpicture}
\begin{axis}[
  width=0.92\linewidth, height=5.4cm,
  ybar, bar width=9pt,
  ymin=-22, ymax=42,
  ylabel={$\Delta$ mean delay vs comparator (\%)},
  ylabel style={font=\sffamily\small},
  symbolic x coords={Old Street, Bloomsbury, Bloomsbury cal., Euston, Euston cal.},
  xtick=data, x tick label style={font=\sffamily\scriptsize},
  ytick={-10,0,10,20,30,40}, yticklabel style={font=\sffamily\scriptsize},
  legend style={font=\sffamily\scriptsize, at={(0.5,1.16)}, anchor=north, legend columns=2, draw=nbgrey},
  axis line style={nbgrey}, tick style={nbgrey},
  ymajorgrids, grid style={nbgrey!25},
  enlarge x limits=0.13,
  nodes near coords, every node near coord/.append style={font=\sffamily\tiny, /pgf/number format/fixed, /pgf/number format/precision=1},
]
\addplot[fill=nbteal!70, draw=nbteal,
  nodes near coords style={xshift=-4pt}] coordinates
  {(Old Street,-11.6) (Bloomsbury,0.8) (Bloomsbury cal.,-8.8) (Euston,15.5) (Euston cal.,8.5)};
\addplot[fill=nbred!60, draw=nbred,
  nodes near coords style={xshift=4pt}] coordinates
  {(Old Street,-9.7) (Bloomsbury,2.8) (Bloomsbury cal.,11.9) (Euston,7.6) (Euston cal.,34.8)};
\legend{fine-tuned, stock}
\draw[nbgrey, dashed] (axis cs:Old Street,0) -- (axis cs:Euston cal.,0);
\end{axis}
\end{tikzpicture}
\caption{Closed-loop delay against the fixed-time floor, $n{=}30$
training-disjoint seeds per cell (negative is better). The distillation
effect is the gap between the bars: largest on the calibrated Bloomsbury
grid, not reproduced at that size on the calibrated corridor.}
\label{fig:headline}
\end{figure*}


\finding{The gridlock was masking the effect, not creating it}
(Table~\ref{tab:calib}, Figure~\ref{fig:headline}). On the
calibrated network the stock controller degrades mean delay by
\risk{11.9\% against the fixed-time floor} and \risk{22.7\% against
MaxPressure}, an order of magnitude beyond the $+2.8\%$ it costs
oversaturated, while both fine-tuned students \finding{beat the fixed-time
floor by 8.8\%}, within $0.03\%$ of MaxPressure. The head-to-head is the
largest, most certain effect in this report: \finding{$18.3\%$} less delay
than the identically-compiled stock model ($[-19.8, -16.8]$,
$p{=}2.0{\times}10^{-21}$), \finding{winning 30 of 30 seeds}, the Phi
student reproducing it ($-18.3\%$, 30/30, $p{=}8.2{\times}10^{-22}$). The
thirty seeds vary SUMO's microscopic noise over one fixed 40\% demand
draw, so they replicate within a scenario rather than across demand; the
mechanism is the one \S\ref{sec:dis-decoupling} develops. Throughput
corroborates it (stock 76.4\% against the fine-tuned 87.4\%), and the
stock arm shows no serving pathology (484 valid proposals of 546
decisions, no fallbacks, normal latency), so the degradation is in the
decisions, not the endpoint.

\finding{This also bounds the scale-saturation result instead of leaving
it a non-rejection}: the 3.8B and 0.6B students differ by $+0.03\%$ with a
95\% CI of $[-0.32, +0.38]$, so any difference between the scales is
within $\pm0.38\%$ of delay. That stops short of demonstrated equivalence,
which would need a margin pre-registered against an operational threshold
(a TOST procedure) this study did not define.

\finding{An oversaturated test network can hide a large effect}: the
original Bloomsbury cells alone would have concluded a marginal $1.8\%$ at
$p{=}.022$. On Euston the same correction did not reveal a comparable
effect, so this is a demonstrated failure mode on one network of two, not
a systematic bias whose size I can quantify. The literature evaluates
predominantly on synthetic grids under demand rarely reported in clearance
terms (\S\ref{sec:bg-rl}), so saturation-driven compression of published
effect sizes deserves attention; reporting completion and teleport counts
alongside delay would expose the regime, though not the compression, since
grid3$\times$3 is oversaturated at 37.8\% completion and still shows
$-19.1\%$.
